\documentclass{article}
\usepackage[utf8]{inputenc}
\newcommand{\comment}[1]{}
\usepackage{tabularx}
\usepackage{hyperref}
\usepackage{mathtools}
\usepackage{float}
\usepackage{subcaption}
\usepackage{textcomp}
\usepackage[a4paper, margin={1in}]{geometry}
\usepackage[english]{babel}
\usepackage{graphicx}
\usepackage{amsmath}
\usepackage{amsfonts}
\usepackage{fontspec}
\setmonofont{FreeMono}
\usepackage{minted}
\setminted{breaklines=true}
\newmintinline[lean]{lean4}{bgcolor=white}
\newminted[leancode]{lean4}{fontsize=\footnotesize}
\usemintedstyle{tango}

\newcolumntype{Y}{>{\centering\arraybackslash}X}

\usepackage{amsthm}
\usepackage{authblk}
\usepackage{biblatex}
\usepackage{mathpartir}
\usepackage{stmaryrd}
\usepackage{amssymb}
\newtheorem{lemma}{Lemma}
\newtheorem{theorem}{Theorem}
\addbibresource{./refs.bib}

\newcommand{\set}{\text{Set}}
\newcommand{\sort}[1]{\text{Sort}\ #1}
\newcommand{\affarr}{\rightarrowtail}
\newcommand{\affetaarr}{\rightarrowtail_{\eta}}
\newcommand{\etaarr}{\rightarrow_{\eta}}

\title{TODO TODO}
\author[1]{Lexi Aiello}
\author[2]{Eric Klavins}
\affil[1]{Department of Philosophy, University of Washington}
\affil[2]{Electrical \& Computer Engineering Department, University of Washington}

\begin{document}
\maketitle
\begin{abstract}
  TODO
\end{abstract}

\section{Semi-Affine Types}

\subsection{SK Combinators}

The SK combinators are the prototypical combinatory logic. They are defined by two reduction rules \eqref{sk_red}.

\begin{split}\label{sk_red}
  K x y &= x \\
  S x y z = x z (y z)
\end{split}

The combinators have been interpreted dependently \cite{altenkirch} \eqref{sk_types}.

\begin{split}\label{sk_types}
  S : \Pi (A : \set) (B : A \rightarrow \set) (C : \Pi (x : A) (y : B x), \set) (x : \Pi (x : A) (y : B x), C x y) (y : \Pi (x : A), B x) (z : A), \gamma z (y z)
  K : \Pi (A : \set) (B : A \rightarrow \set) (x : A) (y : B x), A
\end{split}

Notably, each term in the context is used at most once, except in positive occurrences. Except for the ``final'' codomain of \(S\), each argument type consumes the context affinely. Hence, its types are ``semi-affine.''

\subsection{BCKW Combinators}

The BCKW combinators have the rewrite equations:

\begin{split}
  B x y z &= x (y z)\\
  C x y z &= x z y\\
  K x y &= x\\
  W x y &= x y y
\end{split}

The BCKW combinators are turing-complete. They can be extended with dependent types \eqref{bckw_types}.

\begin{split}\label{bckw_types}
  &B : \Pi (A : \set) (B : \set) (C : B \rightarrow \set) (x : \Pi (x : B), C x) (y : A \rightarrow B) (z : A), C (y x) \\
  &C : \Pi (A : \set) (B : \set) (C : A \rightarrow B \rightarrow \set) (x : \Pi (x : A) (y : B), C x y) (y : B) (z : A), C z y \\
  &K : \Pi (A : \set) (B : A \rightarrow \set) (x : A) (y : B x), A \\
  &W : \Pi (A : \set) (B : A \rightarrow A \rightarrow \set) (x : \Pi (x : A) (y : A), B x y) (y : A), W B y
\end{split}

Similarly to \(SK\), terms in the context are used non-affinely, but only positively.

\section{The Categorical Combinators}

The observation that \(SK\) has semi-affine types is unsatisfying. \(SK\) are well-known to be equivalent to \(\lambda\)-calculus, but there is an arguably closer relative to the \(\lambda\)-calculus: the categorical combinators.

The categorical combinators \autocite{CURIEN1986188} are the basis of the Categorical Abstract Machine \autocite{cmachine}. They are a potential compilation target for the \(\lambda\)-calculus \autocite{cmachine}. Categorical combinators simulate variable substitution through projections over nested tuples (``couples'', or binary products) constituting an ``environment''. Binders are simulated through a curry combinator (\(\Lambda\)) which applies its two arguments as a couple to the body (\(\Lambda(x)\)). The categorical semantics of the calculus are given by cartesian-closed categories \autocite[p. 178]{cmachine}.

\begin{figure}[H]
\begin{center}
\begin{tabular}{l l l}\label{red_untyped_cat}
&(ass) & \((x \circ y) z = x(y z)\) \\
&(fst) & \(Fst(a, b) = a\) \\
&(snd) & \(Snd(a, b) = b\) \\
&(dpair) & \(\langle x, y\rangle z = (x z, y z)\) \\
&(curr) & \(\Lambda(x)yz = x(y, z)\) \\
&(app) & \(App(f, x) = f x\) \\
&(quote) & \(`(x) y = x\)
\end{tabular}
\end{center}
\caption{Reduction rules of the untyped combinators as seen in the Categorical Abstract Machine.}
\end{figure}

Obviously, the only rule that duplicates any terms in scope is \textit{dpair}. Serendipitously, its dependent type can be expressed affinely. The size of its type (and all other types) strictly decrease (unlike in \(SK\) or \(BCKW\)). Our primary contribution in this paper is a technique for normalization in dependent type theories which avoids appealing to denotational semantics based on this insight: by restricting rewrite rules to the affine fragment of the calculus, decidability of \(\beta_{=}\) is automatic. By removing just one combinator (\textit{dpair}), types become strongly normalizing, and type-checking becomes decidable (almost for free). The rest of this paper is dedicated to a practical formalization of the dependently-typed categorical combinators, demonstrating decidability of type-checking, achieved in principal by affine restriction.

% As a result of Sigma types, we can avoid duplication

% Introduction

% SK combinators

% Categorical combinators

% Formalization of dependently-typed SK has yet to be achieved

% Semi-affine types of SK

% Fully affine types of categorical combinators

  % Frame this as a discovery, not a construction

% Affine restriction as a formalization technique

  % Compare and contrast to NbE

% How we used the affine restriction to achieve decidable type-checking for the categorical combinators

% Affine restriction proofs:

  % Decreasing measure

  % Briefly skim typeclass automation

  % \beta-reduction decidability

% Type inference

  % Don't get bogged down into technicals

\end{document}
